\documentclass[10pt,a4paper]{amsart}
\usepackage[margin=19mm]{geometry}
\usepackage{amsmath,amssymb,mathtools}
\usepackage[scr=rsfs,cal=euler]{mathalfa}
\usepackage{xfrac}
\usepackage[unicode,hidelinks,bookmarks=false]{hyperref}
\newcommand{\ie}{i.e.\ }
\newcommand{\cf}{cf.\ }
\newcommand{\resp}{respectively\ }
\newcommand{\Subsection}{\subsection}
\providecommand{\nobreakdash}{\nobreak\hbox{-}}
\setlength{\emergencystretch}{2em}
\multlinegap0pt
\usepackage{amssymb}
\usepackage{xcolor}
\usepackage{tikz}
\newtheorem{proposition}{Proposition}
\title{Metric-Independent Expansiveness}
\author{Alfonso Artigue, Luis Ferrari}
\date{Source: arXiv:2603.20978v1 (CC BY 4.0)}
\begin{document}
\maketitle

\noindent\emph{Source and licence.} Metric-Independent Expansiveness. Version arXiv:2603.20978v1 (CC BY 4.0), distributed by arXiv under the Creative Commons Attribution 4.0 International licence, http://creativecommons.org/licenses/by/4.0/, as recorded in the arXiv licence metadata for this exact version and shown on the abstract page https://arxiv.org/abs/2603.20978v1. Full licence notice: \texttt{licenses/arxiv-2603.20978-CC-BY-4.0.txt}.\par\smallskip

\noindent\textit{Editorial scope.} Counted unit: the article's proposition th:herencia (source0.tex lines 372--374). The article's complete original proof of it is delivered with it (source0.tex lines 376--423, one \texttt{proof} environment of 48 lines). No mathematics is written, repaired or re-derived: the statement and the proof are the article's own lines, in the article's own notation, and no retained line is substituted or re-flowed.  No further source-local context is needed: every cross-reference of every retained line resolves inside the two delivered spans.  Nothing was omitted from any retained span, so the package carries no omission record.  No retained line contains a citation, so the wrapper carries no bibliography and no bibliography entry.  No retained line contains a figure environment, an image-inclusion command or drawing code, so no image is delivered and the package declares no asset dependency.  Editorial formatting only, in addition: an AMS \texttt{amsart} preamble that restates the article's own declarations and macros used by the retained lines, namely the environments \texttt{proposition} on separate counters, together with the standard packages those lines need.  The retained environments print in this document as Proposition 1 in order of appearance, whereas the article prints the same items with the numbering of its own full document.  The complete native source of the article as distributed in the arXiv e-print for this exact version is bundled unchanged as \texttt{sources/196904/source0.tex}.
\begin{proposition}[Inheritance to syndetic subgroups] \label{th:herencia}
Let \( \varphi : G \times X \to X \) be a continuous and cocompactly expansive action and let \( H \) be a syndetic subgroup of \( G \). Then the restricted action \( \varphi|_{H \times X} : H \times X \to X \) is also continuous and cocompactly expansive.
\end{proposition}

\begin{proof}
Let \( K \) be the compact set from the cocompactness of \( \varphi \), \( \mathcal{U} = \{ U_1, \ldots, U_N \} \) be the expansivity cover, and \( F = \{ f_1, \ldots, f_m \} \subset G \) be the finite set such that \( G = F H \).
Each \( f_j^{-1}.K \) is compact, so \( K' := \bigcup_{j=1}^m f_j^{-1}.K \) is compact as well.
For any \( x \in X \), choose \( g \in G \), \( k \in K \) such that \( x = g.k \) (since \( G \cdot K = X \)).
Writing \( g = f_j h \) with \( f_j \in F \), \( h \in H \), we obtain
\[
x = f_j.(h.k) \in H \cdot K'.
\]
Hence, \( H \cdot K' = X \).
Next we define
\[
\mathcal{W}
:= \left\{
  \bigcap_{f \in F_0} f^{-1}.U_{i(f)} \;\middle|\;
  \varnothing \ne F_0 \subseteq F,\;
  i(f) \in \{1, \ldots, N\}
\right\}.
\]
Since \( \mathcal{U} \) and \( F \) are finite, the family \( \mathcal{W} \) is also finite.
Each element of \( \mathcal{W} \) is open (being a finite intersection of open sets) and contains all the sets \( f^{-1}.U_i \), so it covers \( X \).
Adding \( X \setminus K' \), we obtain a finite open cover:
$
\mathcal{W} \cup \{ X \setminus K' \}.
$
Let \( x \ne y \in X \). Since \( \varphi \) is cocompactly expansive, there exists \( g \in G \) such that:
\[
\{ g.x, g.y \} \cap K \ne \varnothing,
\quad \text{and} \quad
\forall U \in \mathcal{U},\;
\{ g.x, g.y \} \not\subset U.
\]
We write \( g = f_j h \) with \( f_j \in F \), \( h \in H \), and define \( x' := h.x \), \( y' := h.y \).
Then \( \{ g.x, g.y \} = f_j.\{ x', y' \} \), and one of these points lies in \( K \), so its preimage lies in \( f_j^{-1}.K \subset K' \).
Thus,
$
\{ x', y' \} \cap K' \ne \varnothing.
$
If \( \{ x', y' \} \subset X \setminus K' \), then applying \( f_j \) gives
\( \{ g.x, g.y \} \subset X \setminus K \), contradicting the assumption.

Now take \( W = \bigcap_{f \in F_0} f^{-1}.U_{i(f)} \in \mathcal{W} \) and fix \( f_k \in F_0 \). Then \( W \subset f_k^{-1}.U_{i(f_k)} \).
If \( \{ x', y' \} \subset W \), then applying \( f_k \) yields:
$
\{ f_k.x', f_k.y' \} \subset U_{i(f_k)} \subset \mathcal{U};
$
which contradicts the expansivity assumption.
Therefore, \( \{ x', y' \} \not\subset W \), and \( \varphi|_{H \times X} \) is cocompactly expansive.
\end{proof}

\end{document}
